\documentclass[11pt]{article}

% ---------------------------------------------------------------------
% Portable package set only -- no arXiv-only class files. Compiles with
% a standard TeX Live / MiKTeX install via pdflatex + bibtex, or latexmk.
% ---------------------------------------------------------------------
\usepackage[utf8]{inputenc}
\usepackage[T1]{fontenc}
\usepackage[margin=1in]{geometry}
\usepackage{amsmath}
\usepackage{amssymb}
\usepackage{amsthm}
\usepackage{graphicx}
\usepackage{booktabs}
\usepackage{array}
\usepackage{multirow}
\usepackage{longtable}
\usepackage{caption}
\usepackage{xcolor}
\usepackage{enumitem}
\usepackage{natbib}
\usepackage{hyperref}

\hypersetup{
  colorlinks=true,
  linkcolor=blue!50!black,
  citecolor=blue!50!black,
  urlcolor=blue!50!black,
  pdftitle={Parametric Inductive Bias vs. Free Feature Learning in High-Frequency Queue Depletion Forecasting},
  pdfauthor={}
}

\newcommand{\todo}[1]{\textcolor{red}{\textbf{[TODO: #1]}}}

\newcommand{\PRAUC}{\mathrm{PR\text{-}AUC}}

\title{Parametric Inductive Bias vs.\ Free Feature Learning:\\
Evaluating a Multivariate Hawkes Process Against Gradient-Boosted\\
Trees for L1 Order-Book Queue Depletion Forecasting}

\author{%
  \todo{Author name(s) and affiliation(s)}
}

\date{\today}

\begin{document}

\maketitle

\begin{abstract}
We ask whether a parametric, unsupervised, maximum-likelihood-constrained
multivariate Hawkes self-exciting point process ($M_2$) generalizes better
than an unconstrained gradient-boosted-tree model ($M_1$) at forecasting
short-horizon L1 limit-order-book queue depletion, given that $M_2$'s
intensity is, by construction, a linear projection of $M_1$'s own
multi-scale exponentially-weighted feature bank and therefore cannot add
information the tree cannot already represent. Using 38 days
(2024-02-22 to 2024-03-30) of unaggregated Binance BTCUSDT perpetual
futures \texttt{trades} and \texttt{bookTicker} data ($1{,}053{,}434{,}969$
causally aligned, typed events across a 10-event microstructure taxonomy),
we find that $M_1$ dominates $M_2$ at every training-sample budget from
2 hours to 31 days on a fixed out-of-sample test week, with no crossover
sample size $N^{*}$: the gap is $0.012$--$0.031$ PR-AUC in $M_1$'s favor and
remains a full order of magnitude above the pre-registered
$\Delta\PRAUC \ge 0.005$ threshold even under a calibration budget scaled
up to 50$\times$ the originally used sample (rolling Giacomini--White test,
$p \approx 0$, statistic $-38.59$). A third model, $M_3$ ($M_1$ plus an
hourly Hawkes-derived spectral-radius regime indicator and linearized
trade-size marks), shows no PR-AUC advantage over $M_1$ on the
500\,ms-horizon Primary Endpoint, but a formal rolling
Giacomini--White test reveals a small, statistically significant
\emph{calibration} (log-loss) improvement over $M_1$ ($p \approx 5.1\times
10^{-12}$; regime-feature-only variant, $p \approx 2.3\times10^{-15}$), an
effect an ablation traces entirely to the regime feature, not the trade
marks. Across an 18-specification secondary grid (2 targets $\times$ 3
depletion thresholds $\times$ 3 horizons), Romano--Wolf FWER-controlled
testing under the full 31-day rolling protocol confirms $M_1 > M_2$ at
every specification (adjusted $p \le 0.0215$), and the same $M_1>M_2$
finding replicates on an independent held-out test week within the
primary window (six of seven budgets tested, gap $0.0107$--$0.0313$
PR-AUC throughout). A cross-asset transfer study extended to six
target assets (BTC $\to$ ETH/SOL/DOGE, then BNB/XRP/ADA) shows the
BTC-fitted branching-ratio matrix $\boldsymbol{\Gamma}$ retains
$78.97\%$--$89.3\%$ of each target asset's own $M_0\!\to\!M_1$
predictive gap when transferred with only the baseline intensity
$\boldsymbol{\mu}(t)$ re-estimated locally; the three assets added under
this robustness extension cluster together just below the floor set by
the original three, a pattern not resolved by the six rungs tested here.
Taken together, the results falsify the hypothesis that MLE-constrained
parametric compression generalizes better than free tree learning at any
tested sample size for ranking-quality forecasting, while identifying a
narrow, calibration-only channel through which a Hawkes-derived regime
signal adds real, if modest, value.
\end{abstract}

\section{Introduction}
\label{sec:introduction}

High-frequency limit-order-book (LOB) dynamics have long been modeled with
two largely separate toolkits. The econometrics tradition favors
parametric, likelihood-based point-process models -- most prominently the
multivariate Hawkes self-exciting process -- which impose a structured,
interpretable form on order-flow clustering and cross-excitation and are
fit by unsupervised maximum likelihood on the event stream itself. The
machine-learning tradition instead favors unconstrained, supervised
function approximators -- gradient-boosted trees, in particular -- fit
directly against a forecasting target on a rich, hand-engineered feature
bank, with no likelihood story and no structural constraint on how inputs
combine. Both approaches are used, sometimes on the same data, to answer
essentially the same applied question: how likely is a given order-book
queue to deplete in the next tens to hundreds of milliseconds?

This paper asks a specific, falsifiable version of the comparison: under
what sample-size regimes, if any, does a maximum-likelihood-constrained
multivariate Hawkes process outperform an unconstrained LightGBM model at
forecasting L1 queue depletion on Binance BTCUSDT perpetual futures, when
both models are built from the \emph{same} underlying multi-scale
exponentially-weighted-moving-average (EWMA) event-count feature bank?

\subsection{The real question is generalization under constraint, not information}
\label{sec:intro-structural-note}

The comparison is deliberately not framed as ``does the Hawkes process see
something the tree cannot.'' It structurally cannot. With the same decay
grid $\beta_p \in \{0.2, 1.0, 5.0, 25.0, 100.0\}\ \mathrm{s}^{-1}$ used by
both models, the fitted multivariate Hawkes intensity
\[
  \boldsymbol{\lambda}(t) = \boldsymbol{\mu}(t) + \sum_{p=1}^{5} \mathbf{A}^{(p)} \mathbf{Z}^{(\beta_p)}(t)
\]
is, by construction, a \emph{linear} function of exactly the same
multi-scale event-count feature bank $\mathbf{Z}^{(\beta_p)}(t)$ that the
gradient-boosted-tree model ($M_1$, Section~\ref{sec:methodology}) is free
to combine nonlinearly, interact, and split on without restriction. A
model whose feature space is a strict linear subspace of another model's
own input space cannot, in the general case, exceed that other model on
any supervised objective by exploiting genuinely new information -- at
best it can match it. The interesting question this comparison can
actually answer is therefore not an information question but a
\emph{generalization} question: does the regularization implicit in an
unsupervised, likelihood-constrained parametric fit -- effectively, a
structured dimensionality reduction of the same feature space down to a
handful of branching-ratio parameters per event-type pair -- generalize
better than letting a flexible nonparametric learner fit that same space
directly, particularly in small-sample, regime-shifted, or
cross-instrument-transfer conditions where the tree has less to learn
from and a strong parametric prior might be expected to help most?

\subsection{Falsification-first framing}
\label{sec:intro-falsification}

This project is explicitly designed to be falsification-first rather than
confirmation-seeking. A single Primary Confirmatory Endpoint is
pre-registered before any result is examined: on the 500\,ms-horizon,
50\%-of-spread adverse-price-shift target, a treatment model must beat
$M_1$ by $\Delta\PRAUC \ge +0.005$ at $\alpha = 0.05$ under a formal,
walk-forward significance test (the rolling Giacomini--White protocol of
Section~\ref{sec:methodology}) for the parametric-compression hypothesis
to be considered supported. Every other comparison in this paper --
the full sample-size learning curve, the 18-specification secondary grid,
the cross-asset transfer ladder, the operational lead-time metric -- is
reported as exploratory evidence, not as a confirmatory test, and is
labeled as such throughout. Where results were surprising or ran counter
to a standing project hypothesis -- most notably, a documented
expectation that untuned LightGBM hyperparameters were flattering $M_1$,
later tested and found backwards (Section~\ref{sec:results-tuning}) -- the
finding is reported plainly rather than smoothed over, consistent with
the project's own stated philosophy that the goal is not to prove the
Hawkes process adds value but to test, honestly, whether it does.

\subsection{Why this matters}
\label{sec:intro-why}

Two practical audiences care about the answer. First, for real-time
microstructure forecasting infrastructure, the choice between a
parametric point-process model and a free tree model is a choice about
sample-efficiency, interpretability, and cross-instrument transferability
versus raw predictive power -- a Hawkes process needs far fewer bytes to
describe, is cheaper to fit incrementally, and, as
Section~\ref{sec:results-transfer} shows, transfers its fitted structure
across instruments in a way a black-box tree model has no native
mechanism to do at all. If the parametric model's generalization
advantage in the small-sample or cross-instrument regime is real, that
trade-off favors the parametric approach in exactly the conditions where
a new instrument or a cold-start deployment has little training data of
its own. Second, for the broader methodological question of when
structured, likelihood-based inductive bias helps versus hurts relative
to unconstrained supervised learning on the same feature space, L1 queue
depletion is a clean, data-rich testbed: the ground-truth event stream is
large (in excess of one billion typed events across the primary window),
the two models share an identical feature basis by design, and the
forecasting target is well defined and mechanically close to the
underlying point process. A clear empirical answer in this setting speaks
to a question broader than any one exchange or instrument.

The remainder of the paper is organized as follows.
Section~\ref{sec:related-work} situates this work relative to prior
literature on Hawkes processes in market microstructure, LOB
queue-depletion modeling, and machine-learning approaches to
high-frequency forecasting. Section~\ref{sec:data} describes the data
pipeline, the 38-day primary window, the 10-event microstructure
taxonomy, and known data-quality constraints. Section~\ref{sec:methodology}
gives the full model ladder ($M_0$--$M_3$), the Hawkes MLE formulation,
the LightGBM setup, and the statistical testing protocols.
Section~\ref{sec:results} reports the empirical results: the fixed-test
learning curve, the rolling Giacomini--White test, the calibration-budget
robustness sweep, the $M_3$ build and its ablation, the 18-specification
Romano--Wolf grid, the cross-asset transfer ladder, the per-$N$
hyperparameter-tuning finding, and the operational lead-time metric.
Section~\ref{sec:limitations} discusses scope limitations, disclosed
simplifications and their resolutions, and reproducibility caveats.
Section~\ref{sec:conclusion} concludes.

\section{Related Work}
\label{sec:related-work}

This section situates the present study relative to four bodies of prior
work: Hawkes processes in market microstructure, limit-order-book (LOB)
queue-position and queue-depletion modeling, machine-learning approaches
to LOB forecasting, and the small literature that compares parametric
point-process models against machine-learned forecasts directly. A fifth
subsection covers the statistical-testing methodology this paper relies
on. Every citation below was independently verified to exist (arXiv,
publisher, SSRN, or Semantic Scholar listing) before inclusion; where a
bibliographic detail such as venue or peer-review status could not be
confirmed, this is noted explicitly rather than presented as certain.

\subsection{Hawkes processes in market microstructure}
\label{sec:rw-hawkes}

Multivariate Hawkes processes are a standard tool for modeling
self-exciting order-flow clustering in financial markets;
\citet{bacry2015hawkes} survey the framework broadly, including
volatility estimation, market stability, contagion, and full
order-book dynamics, and this paper's 10-event marked multivariate
Hawkes process with exponential kernels and a branching-ratio
subcriticality constraint (Section~\ref{sec:methodology-hawkes}) is a
direct descendant of that framework. \citet{bowsher2007modelling}
introduces ``generalised Hawkes'' models for market events allowing
inhibitory effects and cross-day dependence, an early instance of the
same modeling stance taken here -- treating quote and trade events as a
multivariate point process with a full intensity vector, rather than a
single univariate arrival process. \citet{hewlett2006clustering} applies
a bivariate Hawkes process to buy/sell trade arrivals in FX to predict
imbalance, a narrower ancestor of the 10-event taxonomy of
Section~\ref{sec:data-taxonomy} (workshop presentation, not confirmed as
peer-reviewed). \citet{bormetti2015modelling} uses multivariate Hawkes
factor models to capture cojumps across correlated assets, a precedent
for using a shared or constrained Hawkes structure to transport
excitation information across instruments -- directly relevant to the
cross-asset transfer ladder of Section~\ref{sec:results-transfer}, which
transfers a compressed parametric object ($\boldsymbol\Gamma$) across
BTC, ETH, SOL, and DOGE rather than pooling raw training data.

Two papers bear directly on the near-critical branching-ratio values
observed in this study ($\rho(\boldsymbol\Gamma) \in [0.93, 0.98]$
throughout Table~\ref{tab:learning-curve}).
\citet{hardiman2013critical} finds that price-change Hawkes kernels on
S\&P E-mini data (1998--2011) are close to criticality (branching ratio
$\approx 1$) as a persistent stylized fact, consistent with this paper's
own measured range rather than suggesting a calibration artifact. Set
against that, \citet{filimonov2015apparent} warns that regime mixtures
and poor timestamp quality can produce \emph{spuriously} high estimated
branching ratios -- ``apparent criticality'' where the true value is
substantially lower -- a methodological caution directly relevant to
this paper's own tie-break causality sensitivity finding
(Section~\ref{sec:data-tiebreak}): the large swings in
$\boldsymbol\Gamma$'s Frobenius norm under alternative same-millisecond
tie-break conventions are exactly the kind of calibration fragility
\citeauthor{filimonov2015apparent} describe, which motivated reporting
tie-rule sensitivity explicitly rather than a single point estimate for
$\boldsymbol\Gamma$.

\subsection{Limit-order-book queue-position and queue-depletion modeling}
\label{sec:rw-lob}

\citet{cont2014price} shows that order-flow imbalance (OFI) at the best
quotes has an approximately linear relationship with short-horizon price
changes across a broad cross-section of NYSE stocks -- the direct
ancestor of this paper's $M_0$ baseline feature
(Section~\ref{sec:methodology-ladder}), and consistent with $M_0$ alone
already reaching PR-AUC $\approx 0.63$ on the fixed test week
(Table~\ref{tab:learning-curve}). \citet{gould2016queue} similarly finds
L1 bid/ask queue imbalance to be a strong, near-linear predictor of the
next tick move, particularly for large-tick instruments, supporting this
paper's inclusion of absolute queue sizes and rolling percentile rank in
$M_0$. \citet{huang2015simulating} models the LOB as a Markov
queuing system whose event intensities depend on current queue state --
the closest prior framework to this paper's Target B (pure depth
depletion), though a state-Markovian rather than history-Hawkes
formulation: where \citet{huang2015simulating} conditions intensities on
current depth, $M_2$ here instead conditions on recent \emph{event
history} via self- and cross-excitation, a complementary rather than
competing account of the same phenomenon.
\citet{moallemi2017queue} models the economic value of queue position via
a cancellation intensity increasing in queue size, an antecedent of this
paper's \texttt{*\_qty\_down} cancellation-volume event types
(Section~\ref{sec:data-taxonomy}) (working paper; no confirmed
peer-reviewed publication located).

\subsection{Machine learning for LOB forecasting}
\label{sec:rw-ml}

$M_1$'s design -- an unconstrained gradient-boosted-tree model on a
hand-engineered multi-scale event-count feature bank -- sits within an
established line of ML-for-LOB work. \citet{kercheval2015modelling} is an
early instance using multi-class SVMs on LOB-level attributes to predict
mid-price direction, which $M_1$ updates with a modern nonlinear tree
ensemble. \citet{ntakaris2018benchmark} establishes the widely used
FI-2010 benchmark (10 days, 5 Nasdaq Nordic stocks, $\sim$4M samples,
3-class up/stationary/down labels), against which this paper's own
labeling protocol departs in kind: Target A/B
(Section~\ref{sec:methodology-target}) predicts depletion and
adverse-move \emph{events} specifically, rather than the field's dominant
3-class mid-price-direction convention. \citet{zhang2019deeplob} (DeepLOB)
and \citet{wallbridge2020translob} (TransLOB, arXiv preprint, no
confirmed peer-reviewed venue located) represent the deep-learning and
attention-based ends of the same ``unconstrained representation
learning'' side of this paper's central contrast; this paper uses
LightGBM rather than a CNN/LSTM/Transformer as its free-model baseline, a
defensible and considerably cheaper stand-in at this paper's much finer
L1/event-level (rather than book-snapshot-image) granularity.
\citet{sirignano2019universal} shows a pooled, multi-stock LSTM
outperforms asset-specific models and generalizes to held-out stocks --
evidence for a ``universal,'' data-learnable price-formation mechanism
along a different transfer axis than this paper's own cross-asset
transfer ladder (Section~\ref{sec:results-transfer}): where
\citeauthor{sirignano2019universal} pool raw training data across
instruments into one free model, this paper instead transfers a
compressed parametric object ($\boldsymbol\Gamma$) and re-profiles only
the local baseline intensity and marks.

\subsection{Direct comparisons of parametric point-process and ML forecasts}
\label{sec:rw-comparison}

This is the category where the literature search found the least direct
overlap with the present paper's research question, and that gap is
reported here rather than papered over with a loose fit. No paper was
located that runs the same head-to-head this paper runs: a constrained,
MLE-fit multivariate Hawkes process against an unconstrained tree
ensemble, on identical LOB targets and an identical underlying feature
basis, under a formal walk-forward significance protocol. The closest
matches found are as follows. \citet{shi2022state} builds a \emph{hybrid}
neural-Hawkes model (continuous-time LSTMs modeling per-event-type
intensities plus an event-state interaction mechanism) and reports it
outperforms both pure-stochastic and pure-deep-learning baselines on LOB
event-stream prediction -- the closest published instance of a controlled
Hawkes-vs-ML comparison found, but a fusion study arguing that
\emph{combining} the two architectures wins, not a controlled test of
whether the parametric constraint itself generalizes better, which is
this paper's question. \citet{nitturanantha2026forecasting} reports that
a sum-of-exponentials Hawkes process gives the best forecast among
unspecified ``competing models'' for near-term order-flow-imbalance
forecasting on NSE tick data; whether its competitor set includes
tree-ensemble or deep-learning baselines could not be confirmed from the
publicly available abstract, so this paper does not claim it as a
confirmed ML comparison. \citet{mei2017neural} is the foundational
neural-Hawkes paper underlying the line of work
\citeauthor{shi2022state} build on (not finance-specific); it is relevant
background for why $M_2$ in this paper is deliberately kept fully
parametric (closed-form exponential kernels, explicit MLE, explicit
$\boldsymbol\Gamma$) rather than neuralized -- a neural Hawkes intensity
would no longer isolate the specific falsification-first question this
paper asks (Section~\ref{sec:intro-falsification}), namely whether the
parametric \emph{constraint itself} helps or hurts generalization.

In summary: the literature contains many papers on parametric Hawkes
processes for order flow, many papers on ML/deep learning for LOB
prediction, and a small number of hybrid fusion papers showing that
combining the two architectures outperforms either alone. What is
largely absent is a paper that holds the prediction target and the
free-model architecture fixed and asks, as this paper does, whether a
Hawkes-constrained feature representation regularizes better than letting
a tree see the same raw feature bank directly -- an ablation-style
comparison rather than an architecture-fusion or bake-off study. The
learning-curve crossover search of Section~\ref{sec:results-learning-curve}
appears to be a genuinely underexplored angle on the Hawkes-vs-ML
question relative to the literature surveyed here.

\subsection{Statistical methodology}
\label{sec:rw-methodology}

The rolling Giacomini--White protocol used throughout
(Section~\ref{sec:methodology-gw}) is a direct application of
\citet{giacomini2006tests}, which extends Diebold--Mariano/West-style
predictive-ability testing to settings where forecasting models are
estimated or re-estimated on a rolling basis and potentially
misspecified -- exactly the walk-forward, daily-re-estimation design this
paper's GW tests use (Sections~\ref{sec:results-gw}
and~\ref{sec:results-m3-gw}). The Romano--Wolf stepdown FWER control
across the 18-specification secondary grid
(Section~\ref{sec:results-secondary-grid}) follows
\citet{romano2005stepwise}, which frames stepdown multiple-testing
control explicitly around the data-snooping problem in empirical finance
(many strategies or specifications tested on the same data), together
with the companion methodological paper \citet{romano2005exact}
developing the underlying exact/approximate stepdown procedures; the two
are standard to cite jointly, as is done here.

\section{Data}
\label{sec:data}

\subsection{Instrument, feeds, and the bookTicker discontinuation constraint}
\label{sec:data-feeds}

The primary instrument is \texttt{BTCUSDT} on Binance USD-margined
perpetual futures (\texttt{futures/um/daily/}). Two raw public feeds are
used: unaggregated \texttt{trades} (individual execution ticks, not the
\texttt{aggTrades} feed, which collapses simultaneous fills and would
destroy the clustering signal the Hawkes process is meant to measure) and
\texttt{bookTicker} (L1 best-bid/best-ask updates). Both feeds are ordered
on \texttt{transaction\_time} (\texttt{T}), not \texttt{event\_time}
(\texttt{E}), since the latter includes Binance's own engine push
latency and is not a reliable causal ordering key.

The data window is constrained by a permanent upstream limitation:
Binance stopped publishing \texttt{futures/um/daily/bookTicker/BTCUSDT}
after \textbf{2024-03-30}. This is documented in the
\texttt{binance-public-data} public GitHub repository (Issue \#372, filed
November 2024) and remained open and unresolved as of the most recent
check (August 2026, 20+ months after filing) -- it is treated in this
paper as a permanent, not transient, constraint, and rules out any
attempt to use a more recent window. The \texttt{trades} feed itself
remains continuous and current; only \texttt{bookTicker} is affected.
Separately, Binance's own archive is documented (Issue \#305) to have
shipped BTC/ETH \texttt{bookTicker} daily files \emph{out of sequence}
during Jan--Mar 2024, so no assumption of pre-sorted archive files is
made anywhere in the pipeline; all downstream processing sorts explicitly
by \texttt{(transaction\_time, update\_id)} before any diff-based
computation.

\subsection{The 38-day primary window}
\label{sec:data-window}

An initial full-coverage audit surfaced severe raw ordering disorder in
\texttt{bookTicker} between 2024-02-04 and 2024-02-21 -- on some days in
this range, up to $\sim$40\% of a day's rows exhibited a
\texttt{transaction\_time} decrease relative to the previous row (e.g.\
18{,}728{,}015 of 46{,}396{,}761 rows on 2024-02-12), a scale of disorder
far beyond isolated same-millisecond ties. An independent re-audit
confirmed that raw ordering disorder ends precisely at 2024-02-21 --
2024-02-22 onward shows zero raw \texttt{transaction\_time} decreases.
The primary study window is therefore fixed to
\[
  \textbf{2024-02-22 -- 2024-03-30 (38 calendar days)},
\]
the last available continuous window before the \texttt{bookTicker}
discontinuation and after the resolved disorder period. An independent
re-audit of this exact window (38 days, $875{,}126{,}040$ raw
\texttt{bookTicker} rows) found \emph{zero} adjacent
\texttt{transaction\_time} decreases, zero same-millisecond
\texttt{update\_id} regressions, and zero invalid-day rows across every
one of the 38 days. Two days retain known, small activity gaps --
2024-02-28 (largest interval 17:32:13--18:15:20 UTC, 2{,}587 seconds
inactive) and 2024-03-08 (largest interval 16:03:08--17:09:20 UTC,
3{,}972 seconds inactive) -- totaling roughly two hours ten minutes of
real inactivity across the entire 38-day window; every other day is
100.000\% or near-100.000\% active on a 1-second-bin coverage measure.
The excluded 2024-02-04--02-21 disorder window is discussed further as a
scope limitation in Section~\ref{sec:limitations}.

Raw input over the primary window totals $196{,}875{,}099$
\texttt{trades} rows and $875{,}126{,}040$ \texttt{bookTicker} rows
($1{,}072{,}001{,}139$ raw rows combined). After causal alignment and
classification into the 10-event taxonomy of
Section~\ref{sec:data-taxonomy} -- which collapses redundant,
same-direction consecutive \texttt{bookTicker} updates into single typed
transitions -- the aligned event stream contains
\[
  \mathbf{1{,}053{,}434{,}969}\ \text{typed events}: \quad 196{,}875{,}099\ \text{trades} + 856{,}559{,}870\ \text{emitted quote events},
\]
with zero unclassified quote transitions on every one of the 38 days.
This just over one billion typed events is the number referred to
throughout the paper as the primary window's event volume.

\subsection{The 10-event microstructure taxonomy}
\label{sec:data-taxonomy}

Every raw quote or trade transition is mapped to exactly one of ten
mutually exclusive event types, with residual cancellation volume at
price level $p$ defined as
\[
  V_{\text{cancel}}(p) = \max\!\big(0,\ -\Delta Q(p) - V_{\text{trade}}(p)\big).
\]

\begin{center}
\begin{tabular}{@{}cll@{}}
\toprule
\# & Event type & Definition \\
\midrule
1 & \texttt{bid\_qty\_up}    & bid price unchanged, $\Delta Q_b > 0$ \\
2 & \texttt{bid\_qty\_down}  & bid price unchanged, $V_{\text{cancel}}(P_b) > 0$ \\
3 & \texttt{bid\_price\_up}   & bid price increases (spread narrowing / price improvement) \\
4 & \texttt{bid\_price\_down} & bid price decreases (queue depletion / spread widening) \\
5 & \texttt{ask\_qty\_up}    & ask price unchanged, $\Delta Q_a > 0$ \\
6 & \texttt{ask\_qty\_down}  & ask price unchanged, $V_{\text{cancel}}(P_a) > 0$ \\
7 & \texttt{ask\_price\_up}   & ask price increases (queue depletion / spread widening) \\
8 & \texttt{ask\_price\_down} & ask price decreases (spread narrowing / price improvement) \\
9 & \texttt{trade\_buy}      & aggressive buyer-initiated fill \\
10 & \texttt{trade\_sell}     & aggressive seller-initiated fill \\
\bottomrule
\end{tabular}
\end{center}

Table~\ref{tab:event-counts} gives the exact event-type counts over the
full 38-day aligned stream.

\begin{table}[htbp]
\centering
\caption{Event-type counts, 38-day aligned stream (2024-02-22 -- 2024-03-30).}
\label{tab:event-counts}
\begin{tabular}{@{}lr@{}}
\toprule
Event type & Count \\
\midrule
\texttt{bid\_qty\_up}    & 258{,}877{,}869 \\
\texttt{bid\_qty\_down}  & 151{,}038{,}105 \\
\texttt{bid\_price\_up}   & 10{,}624{,}783 \\
\texttt{bid\_price\_down} & 8{,}142{,}038 \\
\texttt{ask\_qty\_up}    & 258{,}664{,}686 \\
\texttt{ask\_qty\_down}  & 150{,}541{,}519 \\
\texttt{ask\_price\_up}   & 8{,}303{,}725 \\
\texttt{ask\_price\_down} & 10{,}367{,}145 \\
\texttt{trade\_buy}      & 98{,}416{,}666 \\
\texttt{trade\_sell}     & 98{,}458{,}433 \\
\midrule
\textbf{Total} & \textbf{1{,}053{,}434{,}969} \\
\bottomrule
\end{tabular}
\end{table}

\subsection{Timestamp alignment and the tie-break causality sensitivity check}
\label{sec:data-tiebreak}

Sub-millisecond ties within a single feed are broken by a deterministic
micro-jitter; within \texttt{bookTicker}, exact-\texttt{transaction\_time}
ties are further broken by \texttt{update\_id} before falling back to
arrival order. The default cross-feed ordering convention on an
exact-millisecond tie between the two feeds is \texttt{trades\_first}
(trade ordered before a same-millisecond \texttt{bookTicker} update).
Because this convention is an arbitrary modeling choice rather than a
directly observable fact about true event order, every Hawkes cross-excitation
result is subjected to a mandatory causality sensitivity check: the
Primary Endpoint and the fitted branching-ratio matrix
$\boldsymbol{\Gamma}$ are recomputed under the inverted rule,
\texttt{book\_first}, and compared.

On the sample day used for initial data inspection (2024-03-30),
exact-millisecond trade/\texttt{bookTicker} collisions accounted for
7.99\% of all quote timestamps (372{,}815 collisions that day), confirming
the sensitivity check is not a minor edge case. At $N=3$d on the 38-day
window, the initial (subsequently found to be confounded, see below)
sensitivity run found $M_2$ PR-AUC of $0.6805835954$ under
\texttt{book\_first} versus $0.6923842563$ under \texttt{trades\_first} --
a difference of $-0.01180066091$ -- while $M_1$ and $M_0$ were essentially
unaffected ($\Delta M_1 = +0.0008535961$, $\Delta M_0 = 0$). The
branching-ratio matrix $\boldsymbol{\Gamma}$ itself changed far more than
the downstream PR-AUC: $43.67\%$ relative Frobenius-norm change, with
$89.22\%$ of the squared change concentrated in the quote$\to$trade
block -- physically, in the same-millisecond touch-price/trade pairs
where the tie rule directly determines whether a touch-price move or the
paired trade is treated as causally prior.

A follow-up investigation isolated this effect more cleanly by holding
the event multiset exactly fixed under an ordering-only inversion (rather
than re-running the full aligner, which also changes trade-volume
attribution at qty-tie edges and confounds ordering with event-typing
effects). Under this cleaner order-only design, the pure tie-order effect
on $\boldsymbol{\Gamma}$ was found to be \emph{larger}, not smaller, than
the original confounded estimate ($61.9\%$ relative Frobenius change,
unmasked), concentrated almost entirely (95\% of squared change) in an
``ignition'' artifact: trades acquire large, spurious self-exciting
coefficients from touch-price events that previously sat strictly before
them at the same timestamp. Masking same-millisecond touch-price/trade
ties out of the Hawkes calibration (\texttt{mask\_touch\_trade}) reduces
this to an $8.4\%$ residual, itself attributable to a disclosed,
narrower and unmasked class of same-millisecond quantity-only quote/trade
ties (6.35\% of all events) whose effect propagates into the
quote$\to$quote block via ordinary multicollinearity in the joint
per-target regression, rather than a defect in the masking logic itself.
Across every configuration tested, the downstream PR-AUC sensitivity of
the pure order-only, masked comparison shrank to $-0.0002365$, roughly
$50\times$ smaller than the original confounded $-0.0118$ figure and
comparable in magnitude to this project's own documented run-to-run
numerical noise floor (Section~\ref{sec:limitations}). The full
investigation, decision log, and its extension to the cross-asset
transfer setting are summarized further in
Sections~\ref{sec:results-transfer} and~\ref{sec:limitations}.

\subsection{Fixed test-week representativeness}
\label{sec:data-testweek}

The learning-curve test set (Section~\ref{sec:methodology}) is fixed to
the final calendar week of the primary window, 2024-03-24 -- 2024-03-30
(a full Sunday-through-Saturday week, guaranteeing all weekdays, one
weekend, and three full 8-hour funding cycles per day regardless of
anchor point). This choice was checked, not assumed: every daily metric
in the test week (trade row count, quote notional, realized volatility,
\texttt{bookTicker} row count, aligned event count) was compared against
the pre-test-week distribution (2024-02-22 -- 2024-03-23) using a robust
$z$-score, $z = (x - \mathrm{median})/(1.4826\,\mathrm{MAD})$, with an
ex-ante flag threshold of $|z| > 3.5$. No test-week observation was
flagged (\texttt{test\_week\_anomalous=false}); the largest absolute
robust $z$-score observed was $2.424538$, for aligned-event count on
2024-03-30.

\section{Methodology}
\label{sec:methodology}

\subsection{Prediction target}
\label{sec:methodology-target}

The Primary Confirmatory Endpoint is Target A (Adverse Price Shift) at
depletion threshold $X = 50\%$, forecast horizon $h = 500$\,ms, and
buffer $\Delta = 50$\,ms: the label is positive if the best price level
(bid for a long-side alert, ask for a short-side alert) moves adversely
within $[t+\Delta,\ t+\Delta+h]$. The pre-registered effect-size threshold
is $\Delta\PRAUC \ge +0.005$ over $M_1$ at $\alpha = 0.05$, fixed before
any result was examined. A secondary/exploratory grid additionally varies
the depletion threshold $X \in \{30\%, 50\%, 70\%\}$ and horizon
$h \in \{100\text{ms}, 500\text{ms}, 2\text{s}\}$ (buffer fixed at 50\,ms)
for both Target A and Target B (Pure Depth Depletion: best price
constant, queue size drops $>X\%$ within the same window), for a nominal
$2 \times 3 \times 3 = 18$ specifications
(Section~\ref{sec:results-secondary-grid}).

\subsection{Model ladder}
\label{sec:methodology-ladder}

Four models are compared, each nested in the features available to the
next:

\begin{itemize}[leftmargin=2em]
\item \textbf{$M_0$ (Baseline L1 LOB):} level-1 order-flow imbalance
  $\mathrm{OFI}_{L1}$, micro-price basis $(P_{\text{micro}} - P_{\text{mid}})$,
  bid-ask spread, and absolute queue sizes with rolling historical
  percentile rank $\mathrm{Rank}(Q(t)) \in [0,1]$. Omitting the percentile
  rank term makes $M_0$ artificially weak and inflates the apparent gain
  of every model built on top of it, so it is included by design.

\item \textbf{$M_1$ (Multi-Scale EWMA Bank -- unrestricted tree learning):}
  $M_0$ plus 10-event $\times$ 5-scale exponential event counts,
  \[
    Z_j^{(\beta_p)}(t) = \int_0^t e^{-\beta_p (t-s)}\,dN_j(s), \qquad
    \beta_p \in \{0.2,\ 1.0,\ 5.0,\ 25.0,\ 100.0\}\ \mathrm{s}^{-1}, \quad
    j \in \{1,\dots,10\},
  \]
  fed to an unconstrained LightGBM gradient-boosted-tree model.

\item \textbf{$M_2$ (Parametric Hawkes -- inductive-bias benchmark):}
  $M_0$ plus the MLE-fitted multivariate Hawkes intensity
  \[
    \boldsymbol{\lambda}(t) = \boldsymbol{\mu}(t) + \sum_{p=1}^{5} \mathbf{A}^{(p)} \mathbf{Z}^{(\beta_p)}(t),
  \]
  using the \emph{same} $\beta$ grid as $M_1$, fed as an additional
  feature set to the same LightGBM architecture. As noted in
  Section~\ref{sec:intro-structural-note}, $\boldsymbol{\lambda}(t)$ is
  structurally a linear function of $M_1$'s own feature bank; the MLE fit
  is unsupervised (fits event likelihood, never the forecasting target)
  and can only impose a regularization prior over the same space, not add
  new information.

\item \textbf{$M_3$ (Structural Non-Nested Dynamics):} $M_1$ plus a rolling
  hourly spectral radius $\rho(\boldsymbol{\Gamma}_t)$ (warm-started,
  non-overlapping hourly refits, forward-filled onto the 1\,Hz feature
  grid as a causal step function) and linearized trade-size marks applied
  only to \texttt{trade\_buy}/\texttt{trade\_sell} events. Both blocks are
  genuinely non-linear or externally-covariate-derived summaries that
  $M_1$'s raw event-count EWMA bank does not contain in any decayed-linear
  form, so $M_3$ is the only model in the ladder that could, structurally,
  exceed $M_1$ for a reason other than overfitting noise. The regime
  feature is explicitly an hourly-resolution indicator and is not expected
  a priori to carry signal at the 500\,ms Primary Endpoint horizon.
\end{itemize}

\subsection{Hawkes MLE formulation}
\label{sec:methodology-hawkes}

The full marked, multivariate intensity used for MLE calibration is
\[
  \lambda_i(t) = \mu_i(t) + \sum_{j=1}^{10} \sum_{p=1}^{P}
    \Big( \alpha_{ij}^{(p)} Z_j^{(\beta_p)}(t) + \alpha_{ij}^{\prime(p)} Z_{j,\mathrm{mark}}^{(\beta_p)}(t) \Big),
\]
with trade-size marks $m_k = \mathrm{Volume}_k / \mathrm{Median}(\mathrm{Volume})_{[t-1\mathrm{h},\,t]}$
computed causally (strictly excluding event $k$ itself from its own
normalizing median) and entering \emph{additively} rather than as a
bilinear $(1+\gamma m)\alpha$ product, which preserves global concavity
of the log-likelihood. The baseline intensity follows a parametric
diurnal and asymmetric funding-jump model,
\[
  \mu_i(t) = \Big(a_0 + \sum_{k=1}^{8}\big(a_k\cos(\tfrac{2\pi k t_{\mathrm{sec}}}{86400}) + b_k\sin(\tfrac{2\pi k t_{\mathrm{sec}}}{86400})\big)\Big)
  \cdot \exp\Big(\sum_{f=1}^{3}\big(\delta_{\mathrm{pre}}\mathbb{1}_{[-5\mathrm{m},0)} + \delta_{\mathrm{post}}\mathbb{1}_{[0,+5\mathrm{m}]}\big)\Big),
\]
with funding times $t_{\mathrm{fund}} \in \{00{:}00, 08{:}00, 16{:}00\}$
UTC and asymmetric pre-/post-funding jump coefficients.

The \emph{dimensionless branching-ratio matrix}
\[
  \boldsymbol{\Gamma} = \sum_{p=1}^{5} \mathbf{A}^{(p)} / \beta_p
\]
governs long-run stability: the process is stationary (subcritical) iff
the spectral radius $\rho(\boldsymbol{\Gamma}) < 1$. The specification
calls for enforcing this constraint \emph{during} MLE; the implementation
instead runs an unconstrained per-target L-BFGS-B fit and applies a
post-hoc uniform radial rescaling of all $\alpha$ (and $\alpha'$)
coefficients to $\rho = 0.995$ only if the unconstrained fit exceeds that
ceiling. This deviation from the specification was tested empirically
across 185 fits spanning the full static $N$-grid, both tie rules, and a
purpose-built hourly stress test, and the projection never activated in
that scan (closest approach $\rho = 0.9875$, 0.75\% headroom); it was
later found to activate exactly once, at $N=31$d, in one hourly window
that coincides with the 2024-03-08 \texttt{bookTicker} gap (a known,
already-documented data artifact, not a general large-$N$ property),
affecting 0.116\% of that budget's training rows -- discussed further in
Section~\ref{sec:limitations}.

MLE calibration subsamples a stratified selection of events per fitting
window rather than using the full raw event stream, for L-BFGS-B
wall-clock tractability; the calibration budget and its sensitivity to
sample-size scaling are examined directly in
Section~\ref{sec:results-calibration}.

\subsection{LightGBM setup}
\label{sec:methodology-lgbm}

$M_0$, $M_1$, $M_2$, and $M_3$ are all realized as LightGBM
gradient-boosted-tree classifiers on their respective feature sets. The
headline learning-curve results
(Section~\ref{sec:results-learning-curve}) use a coarse, fixed
3-bucket hyperparameter heuristic applied identically to $M_1$ and $M_2$
at every training-sample budget $N$. A dedicated follow-up
(Section~\ref{sec:results-tuning}) instead tunes \texttt{num\_leaves},
\texttt{min\_data\_in\_leaf} (equivalently \texttt{feature\_fraction} and
\texttt{lambda\_l2} in the tuned search), and related regularization
parameters independently per $N$ and independently for $M_1$ and $M_2$,
via a causal (time-ordered, no shuffling) train/validation split and
random search, specifically to test whether the untuned heuristic was
biasing the $M_1$-vs-$M_2$ comparison in either direction.

\subsection{Evaluation metrics}
\label{sec:methodology-metrics}

Three metrics are reported throughout: area under the precision-recall
curve (PR-AUC, the primary ranking-quality metric, appropriate given
class-imbalanced positive rates), the Brier score (mean squared error
between predicted probability and binary outcome, a calibration-sensitive
metric), and log-loss (binary cross-entropy, also calibration-sensitive
and the metric on which the formal Giacomini--White test is defined).
PR-AUC and Brier/log-loss can and do disagree in this paper's results
(Section~\ref{sec:results-m3}) -- a model can rank events no better while
still producing better-calibrated probabilities, and the two metrics are
not redundant.

\subsection{Giacomini--White test with Newey--West HAC standard errors}
\label{sec:methodology-gw}

The formal significance test used throughout is the Giacomini--White
\citep{giacomini2006tests} conditional predictive ability (CPA) test, run under a genuinely
out-of-sample rolling walk-forward protocol rather than a single static
train/test split. Concretely: (1) compute the per-event loss differential
$d_t = L(y_t, \hat y_t^{M_a}) - L(y_t, \hat y_t^{M_b})$ (log-loss unless
noted otherwise); (2) aggregate into 1-minute non-overlapping mean blocks
$\bar d_\tau$; (3) run the GW CPA test on $\{\bar d_\tau\}$ with
Newey--West/HAC standard errors; (4) re-estimate both models on a
rolling basis -- retrain every 24 hours on the preceding trailing
$N$-day window, forecast the next day -- so that the test's asymptotics,
which require a genuinely out-of-sample rolling forecast sequence, are
valid. The primary rolling GW protocol used in this paper trains on a
trailing 7-day window, re-estimates daily, and forecasts
2024-02-29 through 2024-03-30 (31 out-of-sample days).

\subsection{Romano--Wolf stepdown FWER control}
\label{sec:methodology-rw}

Family-wise error rate across the 18-specification secondary/exploratory
grid (Section~\ref{sec:results-secondary-grid}) is controlled via a
Romano--Wolf stepdown bootstrap \citep{romano2005stepwise,romano2005exact},
resampled by day-blocks to preserve
intraday dependency structure. The Primary Endpoint
(Section~\ref{sec:methodology-target}) is exempt from this correction --
it is tested once, standalone, at $\alpha = 0.05$, per the pre-registered
protocol. Both a cheaper static 7-day-test-week design (7 day-blocks) and
the full 31-day rolling protocol (31 day-blocks, matching the primary GW
protocol above) were run; results are reported for both, with the 31-day
rolling version treated as the more rigorous of the two.

\subsection{Difference-in-differences placebo control}
\label{sec:methodology-placebo}

A placebo design is specified for validating $M_3$'s structural claim:
permute event-type labels $j \in \{1,\dots,10\}$ at fixed timestamps
(preserving timing, destroying type-specific structure), and require
that $M_3$'s \emph{advantage over} $M_1$ (not merely $M_3$'s own
performance) collapse under scrambling, together with a secondary check
lagging point-process feature inputs by $\tau=10$s and verifying that
predictive advantage decays monotonically toward zero. This placebo
control was not run in the present study (see
Section~\ref{sec:limitations}); given the null PR-AUC result for $M_3$
on the Primary Endpoint (Section~\ref{sec:results-m3}), the placebo test
was judged to have little to falsify at the time and was deprioritized
relative to isolating the source of $M_3$'s calibration effect via
ablation, which was run instead.

\subsection{Operational lead-time metric}
\label{sec:methodology-operational}

Independent of the statistical tests above, an operational metric is
reported at a fixed false-alarm budget (top 1\% alert rate, by each
model's own score distribution): the fraction of adverse events captured
(recall at 1\% alert rate) and the median/mean lead time between an
alert and the qualifying adverse event. No execution simulation,
slippage model, or strategy return is computed -- this remains a
forecasting-quality translation, not a trading study.

\section{Results}
\label{sec:results}

\subsection{Fixed-test learning curve}
\label{sec:results-learning-curve}

The learning-curve experiment fixes the test set to the final 7 calendar
days of the primary window (2024-03-24 -- 2024-03-30, $579{,}593$ rows,
target prevalence $0.3440345208$) and varies the training-sample budget
$N \in \{2\text{h}, 1\text{d}, 3\text{d}, 7\text{d}, 14\text{d}, 23\text{d},
31\text{d}\}$ as rolling backward windows ending at the test-set boundary,
so that every $N$ is compared against identical held-out observations.
Hawkes calibration at every $N$ used a fixed 100{,}000 type-stratified
event subsample and the seasonal-profile baseline; all ten target-row
optimizations converged at every $N$. Table~\ref{tab:learning-curve}
reports the full result.

\begin{table}[htbp]
\centering
\caption{Fixed-test learning curve, PR-AUC and Brier score by training
budget $N$. Test set: 2024-03-24 -- 2024-03-30 (579{,}593 rows,
prevalence 0.3440). $\rho(\boldsymbol\Gamma)$ is the fitted spectral
radius; the post-hoc stationarity projection did not bind at any $N$
(scale $=1.0$ throughout).}
\label{tab:learning-curve}
\resizebox{\textwidth}{!}{%
\begin{tabular}{@{}lrrrrrrrr@{}}
\toprule
$N$ & Train rows & $M_0$ PR-AUC & $M_1$ PR-AUC & $M_2$ PR-AUC & $M_0$ Brier & $M_1$ Brier & $M_2$ Brier & $\rho(\boldsymbol\Gamma)$ \\
\midrule
2h  & 7{,}200     & 0.5705600745 & 0.6869820960 & 0.6559263216 & 0.2223043455 & 0.1916784603 & 0.2098252854 & 0.9382741897 \\
1d  & 82{,}799    & 0.5959138042 & 0.7008203642 & 0.6794294650 & 0.2083216397 & 0.1841098182 & 0.1873945943 & 0.9327609059 \\
3d  & 248{,}397   & 0.6252981828 & 0.7050007357 & 0.6923842563 & 0.1891395145 & 0.1680171744 & 0.1707870065 & 0.9542078309 \\
7d  & 579{,}593   & 0.6276258471 & 0.7058287032 & 0.6918090742 & 0.1858112027 & 0.1653064068 & 0.1685688648 & 0.9549810294 \\
14d & 1{,}159{,}186 & 0.6295886303 & 0.7066682758 & 0.6932490467 & 0.1911223953 & 0.1652326226 & 0.1685409388 & 0.9604496341 \\
23d & 1{,}899{,}085 & 0.6296402443 & 0.7062285858 & 0.6913830866 & 0.2001300809 & 0.1673230436 & 0.1723540581 & 0.9452233127 \\
31d & 2{,}558{,}894 & 0.6282185486 & 0.7058569833 & 0.6921162539 & 0.2139888824 & 0.1710036128 & 0.1756603864 & 0.9561896057 \\
\bottomrule
\end{tabular}%
}
\end{table}

$M_1$ leads $M_2$ on PR-AUC at every $N$, by a margin of roughly
$0.012$--$0.031$ PR-AUC. The gap narrows sharply from $N=2$h to
$N=3$d and then largely flattens from $3$d through $31$d; no crossover
$N^{*}$ at which $M_1$ overtakes $M_2$ is observed at any tested budget
-- if anything $M_2$ never catches up, contrary to a naive
small-sample-favors-the-parametric-model intuition. As discussed in
Section~\ref{sec:results-calibration}, this exact magnitude/shape was
subsequently found to be partly an artifact of a fixed (not $N$-scaled)
Hawkes calibration event budget; the direction of the finding --
$M_1 > M_2$ throughout, no crossover -- was confirmed robust under a
corrected, $N$-scaled calibration budget.

\subsection{Rolling Giacomini--White test}
\label{sec:results-gw}

The rolling protocol (trailing 7-day training window, daily
re-estimation, forecasting 2024-02-29 through 2024-03-30) produced 31
out-of-sample days, $2{,}561{,}477$ OOS rows, and $42{,}694$ one-minute
loss blocks. With the loss differential defined as $M_1$ log-loss minus
$M_2$ log-loss, the mean differential was $-0.01000890845$, the GW
statistic was $-38.58596988$, and the two-sided $p$-value was
computationally $0.0$ (floating-point underflow at the reported
precision). The negative sign favors $M_1$: this is the formal,
walk-forward confirmation of the pattern in
Table~\ref{tab:learning-curve}, at overwhelming statistical significance.

\subsection{Calibration-budget robustness}
\label{sec:results-calibration}

A control investigation established that the 100{,}000-event Hawkes
calibration subsample used throughout Table~\ref{tab:learning-curve} is a
\emph{fixed} constant, not scaled with $N$ -- it represents 8.086\% of
raw calibration events available at $N=2$h but only 0.011\% at $N=31$d,
even as the tree-visible training-row count grows roughly $355\times$
across the same grid. A single control fit at $N=14$d with the budget
raised $10\times$ (to 1{,}000{,}000 events) moved $M_2$'s PR-AUC by
$+0.0012271$ (closing 9.24\% of its gap to $M_1$ at that $N$) -- a real,
above-noise effect, roughly 6--14$\times$ the project's own documented
run-to-run reproducibility noise floor ($\sim 0.0001$--$0.0002$).

A full $N$-scaled sweep was then run, with budgets chosen per $N$
(100{,}000 at 2h, rising to 5{,}000{,}000 at 31d), controlled against a
matched same-machine flat-100{,}000 rerun to isolate the budget effect
from an unrelated cross-machine Hawkes non-bit-reproducibility finding
(Section~\ref{sec:limitations}). Table~\ref{tab:calibration-sweep}
reports the result.

\begin{table}[htbp]
\centering
\caption{$N$-scaled Hawkes calibration budget sweep versus a matched
same-machine flat-100{,}000-event baseline. Gap = $M_1 - M_2$ PR-AUC.
$M_0$/$M_1$ matched exactly between all runs at every $N$ (neither
touches the Hawkes fit).}
\label{tab:calibration-sweep}
\begin{tabular}{@{}lrrrrrr@{}}
\toprule
$N$ & Budget & $M_2$, 100k & $M_2$, scaled & $\Delta M_2$ & Gap, 100k & Gap, scaled (closed) \\
\midrule
2h  & 100{,}000     & 0.6575473 & 0.6575473 & 0.0      & 0.02943 & 0.02943\ (0.0\%) \\
1d  & 300{,}000     & 0.6796216 & 0.6804154 & +0.00079 & 0.02120 & 0.02040\ (3.7\%) \\
3d  & 1{,}000{,}000 & 0.6924211 & 0.6931325 & +0.00071 & 0.01258 & 0.01187\ (5.7\%) \\
7d  & 2{,}000{,}000 & 0.6916795 & 0.6934870 & +0.00181 & 0.01415 & 0.01234\ (12.8\%) \\
14d & 3{,}000{,}000 & 0.6935275 & 0.6947720 & +0.00124 & 0.01314 & 0.01190\ (9.5\%) \\
23d & 4{,}000{,}000 & 0.6913924 & 0.6943620 & +0.00297 & 0.01484 & 0.01187\ (20.0\%) \\
31d & 5{,}000{,}000 & 0.6920127 & 0.6943959 & +0.00238 & 0.01384 & 0.01146\ (17.2\%) \\
\bottomrule
\end{tabular}
\end{table}

Every $N$ shows a real, positive $\Delta M_2$; the gap closes
$3.7\%$--$20.0\%$ depending on $N$, growing roughly with $N$ (consistent
with the fixed budget representing a shrinking fraction of available
data as $N$ grows). \textbf{The qualitative finding is confirmed robust,
not merely ``unlikely to reverse'': even at the most aggressive budget
tested} ($N=31$d, 5{,}000{,}000 events -- still only 0.53\% of the
$936{,}567{,}648$ raw events available at that $N$), \textbf{$M_1$ leads
$M_2$ by 0.01146 PR-AUC, a full order of magnitude above the
pre-registered $\Delta\PRAUC \ge 0.005$ threshold, in the opposite
direction. No crossover $N^{*}$ exists anywhere in the tested range under
either calibration budget.} What Table~\ref{tab:learning-curve} gets
wrong is the \emph{exact magnitude} of the $M_1$--$M_2$ gap at each $N$
(overstated by 4--20\%), not the direction of the headline finding.

\subsection{$M_3$: structural non-nested dynamics}
\label{sec:results-m3}

$M_3$ (Section~\ref{sec:methodology-ladder}) was built and evaluated on
the fixed test week at four budgets, $N \in \{3\text{d}, 7\text{d},
14\text{d}, 31\text{d}\}$. Table~\ref{tab:m3-comparison} reports the full
$M_0$--$M_3$ comparison.

\begin{table}[htbp]
\centering
\caption{$M_0$--$M_3$ comparison, fixed test week. $\Delta(M_3 - M_1)$ is
the raw PR-AUC difference (not a significance test; see
Table~\ref{tab:m3-gw} for the formal test).}
\label{tab:m3-comparison}
\begin{tabular}{@{}lrrrrrr@{}}
\toprule
$N$ & Train rows & $M_0$ & $M_1$ & $M_2$ & $M_3$ & $\Delta(M_3 - M_1)$ \\
\midrule
3d  & 248{,}397     & 0.6252981828 & 0.7050007357 & 0.6922661830 & 0.7048622443 & $-0.0001384914$ \\
7d  & 579{,}593     & 0.6276258471 & 0.7058287032 & 0.6916416162 & 0.7061519538 & $+0.0003232506$ \\
14d & 1{,}159{,}186 & 0.6295886303 & 0.7066682758 & 0.6933895106 & 0.7070705525 & $+0.0004022767$ \\
31d & 2{,}558{,}894 & 0.6282185486 & 0.7058569833 & 0.6921377668 & 0.7063012232 & $+0.0004442399$ \\
\bottomrule
\end{tabular}
\end{table}

\noindent
Brier scores (lower is better) at the same budgets show $M_3$ at or
slightly below $M_1$ from $N=7$d onward (e.g.\ at $N=31$d: $M_1$
$0.1710036128$ vs.\ $M_3$ $0.1700600768$), a hint -- confirmed formally
below -- that $M_3$'s effect, if real, is a calibration effect rather than
a ranking effect. Every $\Delta(M_3-M_1)$ in Table~\ref{tab:m3-comparison}
is at least an order of magnitude below the $+0.005$ pre-registered
threshold, and the sign is inconsistent (negative at $N=3$d, small and
positive from $7$d onward) -- on raw PR-AUC alone, this is a null result
on the Primary Endpoint at every tested $N$, consistent with the
specification's own expectation that the hourly regime indicator is
``not expected to carry 100\,ms-horizon micro-timing signal'' against a
500\,ms-horizon endpoint. ($M_2$'s values in this table differ from
Table~\ref{tab:learning-curve} by $\sim 0.0001$--$0.0002$ at every $N$; this
is the same cross-machine Hawkes non-bit-reproducibility discussed in
Section~\ref{sec:limitations}, not an $M_3$-related effect -- confirmed
because $M_0$/$M_1$, which never touch the Hawkes fit, reproduce their
Table~\ref{tab:learning-curve} values to full displayed precision.)

A control fit at $N=14$d with the Hawkes calibration budget raised
$10\times$ (matching Section~\ref{sec:results-calibration}) moved
$\Delta(M_3-M_1)$ by only $-0.0000191$ -- an order of magnitude
\emph{below} the noise floor that flagged $M_2$'s analogous move as real
-- confirming, for a structural reason (neither of $M_3$'s two extra
feature blocks reads the Hawkes MLE's own fitted intensities; the regime
block uses its own independently-capped 250{,}000-event hourly tracker and
the marks block uses no MLE at all), that $M_3$'s null PR-AUC result is
not a calibration-budget artifact.

\subsubsection{The formal rolling GW test reverses the qualitative reading}
\label{sec:results-m3-gw}

A formal rolling Giacomini--White test for $M_3$ (same 7-day trailing
window, daily re-estimation, 2024-02-29 -- 2024-03-30 protocol as the
official $M_1$-vs-$M_2$ result) was subsequently run, and produced an
unexpected but real result, reported in full in
Table~\ref{tab:m3-gw}. A sanity check first confirmed that $\mathrm{GW}(M_1,M_2)$
recomputed fresh on this same pipeline closely reproduces the official
Section~\ref{sec:results-gw} numbers (mean differential $-0.010035$ vs.\
official $-0.010009$; statistic $-38.617$ vs.\ official $-38.586$, small
deviations consistent with the same cross-machine non-reproducibility
noted above, not a new concern).

\begin{table}[htbp]
\centering
\caption{Formal rolling Giacomini--White test, $M_1$ vs.\ $M_3$, both the
combined regime+marks definition and a regime-feature-only ablation
variant. 31-day rolling protocol, 42{,}694 minute blocks, 2{,}561{,}477
OOS rows. Positive differential means $M_3$ has lower (better) log-loss
than $M_1$.}
\label{tab:m3-gw}
\begin{tabular}{@{}lrr@{}}
\toprule
Quantity & $M_3$ = regime + marks & $M_3$ = regime only \\
\midrule
Mean log-loss differential ($M_1 - M_3$) & $+0.0010446644809617568$ & $+0.0011867723437099665$ \\
GW statistic & $+6.902180900472671$ & $+7.9226506944173085$ \\
$p$-value & $5.121016234527669 \times 10^{-12}$ & $2.325000489836069 \times 10^{-15}$ \\
\bottomrule
\end{tabular}
\end{table}

\textbf{$M_3$ clears the significance bar at $\alpha=0.05$ under the
project's own formal test, on log-loss, at $p \approx 5.1 \times
10^{-12}$}, an order of magnitude beyond the significance threshold --
the opposite qualitative conclusion from the fixed-test-week PR-AUC
reading above. This is not a contradiction once the difference between
metric (log-loss/calibration vs.\ PR-AUC/ranking) and protocol
(31 independent daily walk-forward folds vs.\ one static test week) is
made explicit: a small, consistently one-directional calibration effect
invisible against day-to-day noise in a single week can still reach
overwhelming significance once aggregated over 31 days. Day-by-day
PR-AUC in the rolling run is genuinely split (roughly half the 29
forecast days favor $M_3$, half favor $M_1$, by margins of
$\pm 0.0001$--$0.001$), consistent with a calibration effect rather than a
ranking effect and inconsistent with a leakage or bug explanation for the
GW result. The effect size itself is an order of magnitude smaller than
the $M_1$-vs-$M_2$ effect (mean differential $+0.0010$ here vs.\
$-0.0100$ for $M_1$-vs-$M_2$).

\textbf{The correct combined statement is: $M_3$ clears the formal
significance bar on log-loss, but has not been shown to clear the
pre-registered $\Delta\PRAUC \ge +0.005$ magnitude bar}, which is defined
on a different metric and was already measured to be missed by roughly an
order of magnitude at every tested $N$.

\subsubsection{Ablation: the regime feature, not the trade-size marks, drives the effect}
\label{sec:results-m3-ablation}

A static ablation at $N=14$d and $N=31$d isolated $M_3$'s two extra
feature blocks (Table~\ref{tab:m3-ablation}): $M_1$ plus regime only,
$M_1$ plus marks only, and $M_1$ plus both (the full $M_3$ definition).

\begin{table}[htbp]
\centering
\caption{$M_3$ ablation: regime feature vs.\ trade-size marks, fixed test
week. Values are deltas versus $M_1$.}
\label{tab:m3-ablation}
\begin{tabular}{@{}llrr@{}}
\toprule
$N$ & Variant & $\Delta$PR-AUC vs.\ $M_1$ & $\Delta$Brier vs.\ $M_1$ \\
\midrule
14d & +regime only  & $+0.0002044$ & $-0.0001865$ (better) \\
14d & +marks only   & $+0.0000600$ & $+0.0001075$ (worse) \\
14d & +both ($M_3$) & $+0.0003832$ & $-0.0001490$ (better) \\
31d & +regime only  & $+0.0001652$ & $-0.0010505$ (better) \\
31d & +marks only   & $+0.0001522$ & $+0.0002226$ (worse) \\
31d & +both ($M_3$) & $+0.0003321$ & $-0.0009384$ (better) \\
\bottomrule
\end{tabular}
\end{table}

The pattern is consistent across both $N$ and both metrics: the regime
feature $\rho(\boldsymbol{\Gamma}_t)$ drives essentially all of the
calibration (Brier) improvement, and the marks block makes Brier slightly
\emph{worse} on its own at both $N$; ``both'' is not simply additive --
full $M_3$'s Brier improvement is smaller in magnitude than regime alone,
indicating marks partially dilutes the regime feature's calibration
benefit when combined. A follow-up formal rolling GW test isolating the
regime-only variant (Table~\ref{tab:m3-gw}, right column) confirms this:
dropping marks entirely makes the effect \emph{larger and more
significant}, not smaller (mean differential $+0.00104 \to +0.00119$,
$+14\%$; $p$-value tightening three orders of magnitude, $5.1\times
10^{-12} \to 2.3\times10^{-15}$). \textbf{The hourly $\rho(\boldsymbol{\Gamma}_t)$
regime feature alone produces a small but highly significant log-loss
improvement over $M_1$ -- a cleaner result than the combined $M_3$
definition, because the marks block was mildly diluting it. This still
does not clear the pre-registered PR-AUC magnitude threshold.}

\subsection{18-specification secondary grid and Romano--Wolf FWER control}
\label{sec:results-secondary-grid}

The secondary/exploratory grid (2 targets $\times$ 3 depletion thresholds
$X$ $\times$ 3 horizons $h$) was evaluated under the full 31-day rolling
protocol (trailing 7-day training window, daily re-estimation,
2024-02-29 -- 2024-03-30), with the shared unsupervised Hawkes calibration
computed once per forecast day (not once per specification) since it
never sees the target, and only the 18 $\times$ 2 supervised LightGBM
trees refit daily. A codebase quirk is disclosed here rather than hidden:
\texttt{target\_a}'s implementation does not take a depletion-fraction
argument (only Target B genuinely varies by $X$), so the nominal 18
named specifications contain only 12 distinct tests (3 for Target A, one
per horizon; 9 for Target B). All 18 named columns are nonetheless
reported, as specified. Table~\ref{tab:romano-wolf} gives the full
31-day rolling result; a cheaper static 7-day-test-week first pass (7
day-blocks) found the same qualitative pattern with a looser bound
(adjusted $p \le 0.0375$, vs.\ $\le 0.0215$ under the full rolling
protocol reported here).

\begin{table}[htbp]
\centering
\caption{Romano--Wolf FWER-controlled test across the 18-specification
secondary grid, full 31-day rolling protocol. Mean diff is $M_1 - M_2$
log-loss (negative favors $M_1$). Target A's three $X$-values are
byte-identical per horizon (see text) and shown collapsed.}
\label{tab:romano-wolf}
\begin{tabular}{@{}lrr@{}}
\toprule
Spec & Mean diff ($M_1-M_2$) & Romano--Wolf adjusted $p$ \\
\midrule
target\_a\_x\{30,50,70\}\_h100  & $-0.008972$ & 0.004 \\
target\_a\_x\{30,50,70\}\_h500  & $-0.010056$ & 0.003 \\
target\_a\_x\{30,50,70\}\_h2000 & $-0.009451$ & 0.006 \\
target\_b\_x30\_h100   & $-0.003294$ & 0.0215 \\
target\_b\_x30\_h500   & $-0.006057$ & 0.006 \\
target\_b\_x30\_h2000  & $-0.006895$ & 0.012 \\
target\_b\_x50\_h100   & $-0.004207$ & 0.0215 \\
target\_b\_x50\_h500   & $-0.006574$ & 0.006 \\
target\_b\_x50\_h2000  & $-0.007942$ & 0.006 \\
target\_b\_x70\_h100   & $-0.006245$ & 0.012 \\
target\_b\_x70\_h500   & $-0.008168$ & 0.004 \\
target\_b\_x70\_h2000  & $-0.009420$ & 0.003 \\
\bottomrule
\end{tabular}
\end{table}

\textbf{Every specification shows the same sign ($M_1$ beats $M_2$) and
every Romano--Wolf-adjusted $p$-value clears $\alpha=0.05$}, with a
maximum adjusted $p$ of $0.0215$ under the full rolling protocol. The
$M_1 > M_2$ finding is not an artifact of the single Primary Endpoint
target -- it holds, FWER-controlled, across both target definitions and
all three horizons. Target B's effect sizes are the largest and most
significant of the grid, consistent with Target B (pure depth depletion)
being a more mechanically order-flow-driven target than Target A.

\subsection{Cross-asset transfer ladder}
\label{sec:results-transfer}

The transfer protocol fixes the branching-ratio matrix
$\boldsymbol{\Gamma}$ estimated from a BTC source window
($N=3$d, 2024-03-21 -- 2024-03-24; $\rho(\boldsymbol\Gamma) = 0.9570094762$,
10/10 converged) and re-estimates only the baseline intensity
$\boldsymbol{\mu}(t)$ (via the profiled-MLE estimator described below) on
each target asset, over a matched 7-day window (2024-03-24 -- 2024-03-30,
5-day train / 2-day test split). This was first run on ETH, SOL, and DOGE,
then, as a robustness extension, independently repeated on three further
liquid USDT-margined perpetuals (BNB, XRP, ADA) using an identical
protocol and a freshly-refit BTC source $\boldsymbol\Gamma$
($\rho(\boldsymbol\Gamma)=0.9570094762$, matching the original fit to 10
decimal places). Table~\ref{tab:transfer-synthesis} synthesizes the
result across all six rungs.

\begin{table}[htbp]
\centering
\caption{Cross-asset transfer ladder synthesis, all six rungs tested.
$M_2$ uses BTC's fitted $\boldsymbol\Gamma$ transferred directly
(\texttt{trades\_first} tie treatment) with only $\boldsymbol\mu(t)$
re-estimated locally via the profiled-MLE estimator. ``\% gap captured''
is $(M_2-M_0)/(M_1-M_0)$. ETH/SOL/DOGE are the original rung; BNB/XRP/ADA
are the robustness-extension rung, added under an identical protocol.}
\label{tab:transfer-synthesis}
\begin{tabular}{@{}lrrrrrr@{}}
\toprule
Asset & $M_0$ & $M_1$ & $M_2$ (transferred $\boldsymbol\Gamma$) & gap($M_1-M_0$) & gap($M_2-M_0$) & \% gap captured \\
\midrule
ETH  & 0.4464 & 0.5132 & 0.5009 & 0.0667 & 0.0545 & 81.7\% \\
SOL  & 0.6345 & 0.7008 & 0.6937 & 0.0663 & 0.0592 & 89.3\% \\
DOGE & 0.6694 & 0.7505 & 0.7377 & 0.0811 & 0.0683 & 84.2\% \\
\addlinespace
BNB  & 0.4290 & 0.5278 & 0.5083 & 0.0987 & 0.0793 & 80.27\% \\
XRP  & 0.2378 & 0.3280 & 0.3090 & 0.0903 & 0.0713 & 78.97\% \\
ADA  & 0.2650 & 0.3335 & 0.3197 & 0.0685 & 0.0547 & 79.88\% \\
\bottomrule
\end{tabular}
\end{table}

\textbf{The same qualitative ordering as every BTC result holds on all
six target assets with a $\boldsymbol\Gamma$ matrix fit entirely on BTC}:
$M_0 < M_2(\text{transferred}) < M_1$, on every rung. The transferred BTC
$\boldsymbol\Gamma$ carries real, substantial signal on every asset
tested, capturing $78.97\%$--$89.3\%$ of each asset's own
$M_0\!\to\!M_1$ gap.

\textbf{The original three-rung reading of ``no evidence of degradation''
does not extend cleanly to six rungs, and this is reported as found, not
adjusted to fit the prior framing.} ETH/SOL/DOGE spanned $81.7\%$--$89.3\%$,
with SOL/DOGE if anything capturing \emph{more} of the gap than ETH -- the
opposite of a naive ``transfer degrades with distance from the source
asset'' expectation. The three robustness-extension assets (BNB, XRP, ADA)
instead cluster tightly at $78.97\%$--$80.27\%$ (a $1.3$-point spread,
versus the original three's $7.6$-point spread), \emph{below} the
original range's floor of $81.7\%$. With only six rungs and a single BTC
source window, this is not enough evidence to conclude the ladder
degrades systematically past DOGE -- but it is no longer accurate to
describe the six-rung picture as showing no pattern at all. A concrete,
untested follow-up: whether this two-cluster split correlates with a
structural asset property (relative tick size, quote/trade volume ratio,
maker/taker mix) rather than being coincidental. A related, disclosed
irregularity: XRP and ADA's test-week target prevalence
($0.088673$, $0.117598$) falls \emph{below} BTC's own ($\sim 0.34$),
breaking the pattern set by ETH and SOL/DOGE (all higher-prevalence than
BTC, consistent with the tick-size rationale motivating the ladder's
asset ordering in the first place).

Three design simplifications were tested rather than assumed and found not
to matter for the original ETH/SOL/DOGE rung: a literal daisy-chain design
(each asset's own re-estimated $\boldsymbol\Gamma$ feeding the next rung,
rather than transferring directly from BTC to every target) gave
essentially the same answer as the direct design (SOL/DOGE $M_2$ PR-AUC
deltas of $-0.000774$/$+0.000943$, no consistent direction); a
same-millisecond quantity-tie dual-$\boldsymbol\Gamma$ robustness check
(Section~\ref{sec:data-tiebreak}) found a negligible effect transferred
onto each target asset's own feature distribution (ETH $+0.0002$, SOL
$-0.0009$, DOGE $-0.0004$ PR-AUC, no monotonic trend down the ladder,
not re-checked for BNB/XRP/ADA); and replacing the initial closed-form
empirical-rate $\boldsymbol\mu$ estimator with the properly derived,
fixed-$\boldsymbol\alpha$ profiled-MLE estimator used throughout this
subsection (solved by 1-D bisection, validated to reproduce BTC's own
jointly-fit $\boldsymbol\mu$ to within 0.08\% relative error including
boundary cases) changed every one of the original eight transfer
evaluations' PR-AUC by at most $0.0003$, with no consistent direction --
this closes what had been the last disclosed simplification in the
original transfer-ladder design (Section~\ref{sec:limitations}).

\subsection{Robustness: an independent held-out test week}
\label{sec:results-robustness}

Because BTCUSDT \texttt{bookTicker} data is not published past
2024-03-30 (Section~\ref{sec:data}), a genuinely independent external
calendar window could not be constructed for a further robustness check;
this was re-verified directly against the live archive rather than
assumed, including inspecting a monthly rollup file one step past the
daily cutoff, which turned out to contain only 6h47m of 2024-04-01 data --
not a usable window. As the nearest feasible alternative, the core
$M_0$/$M_1$/$M_2$ learning curve (Section~\ref{sec:results-learning-curve})
was rerun against an independent, non-overlapping, representativeness-checked
7-day test week drawn from earlier in the same 38-day window
(2024-03-17 -- 2024-03-23, immediately preceding the official test week;
no anomaly flags, max $|z|=1.514$), at six of the seven original budgets
($N \in \{2\text{h},1\text{d},3\text{d},7\text{d},14\text{d},23\text{d}\}$;
$31$d was not reachable with only 24 pretest days available before this
earlier week).

\textbf{The $M_1 > M_2$ finding replicates cleanly}: $M_1$ leads $M_2$ at
every one of the six budgets tested, by $0.0107$--$0.0313$ PR-AUC, a full
order of magnitude above the pre-registered threshold throughout, with no
crossover. Absolute PR-AUC levels are not comparable to
Table~\ref{tab:learning-curve} -- this week's target prevalence
($0.5109$) is roughly $1.5\times$ the official week's, which alone
inflates every model's baseline PR-AUC -- so this should be read as
confirming the \emph{direction} of the finding on an independent week,
not as a magnitude-matched replication. One genuine, disclosed
difference: the gap shrinks monotonically with $N$ on this alternate
week ($0.0313$ at 2h down to $0.0107$ at 23d), a cleaner pattern than the
official week's non-monotonic shape (Table~\ref{tab:learning-curve});
with only two weeks compared, whether this reflects real week-to-week
variability in the learning-curve trajectory or is a coincidence of two
point estimates is not resolved here. Because this alternate week is
drawn from the same 38-day archive and market regime as the official
week, it should be read as a within-window robustness check, not
independent-period external validity, which Section~\ref{sec:limitations}
discusses as the paper's most significant remaining open scope
limitation.

\subsection{Per-$N$ LightGBM hyperparameter tuning}
\label{sec:results-tuning}

A standing project hypothesis, carried across several earlier working
sessions, held that $M_1$'s hyperparameters being undertuned relative to
$M_2$'s (both used the same coarse 3-bucket heuristic) understated $M_1$'s
true advantage, and that a proper per-$N$ retune would \emph{widen} the
$M_1$-vs-$M_2$ gap. This was tested directly with an independent causal
random search per $N$ and per model (Section~\ref{sec:methodology-lgbm})
and found backwards. Table~\ref{tab:per-n-tuning} reports the result.

\begin{table}[htbp]
\centering
\caption{Per-$N$ independent LightGBM tuning, $M_1$ and $M_2$, versus the
untuned coarse 3-bucket heuristic. Gap = $M_1 - M_2$ PR-AUC.}
\label{tab:per-n-tuning}
\begin{tabular}{@{}lrrrrrr@{}}
\toprule
$N$ & $M_1$ untuned & $M_1$ tuned & $M_2$ untuned & $M_2$ tuned & gap untuned & gap tuned ($\Delta$) \\
\midrule
2h  & 0.68698 & 0.68315 & 0.65755 & 0.65317 & 0.02943 & 0.02998\ ($+0.00054$) \\
1d  & 0.70082 & 0.70111 & 0.68041 & 0.68015 & 0.02041 & 0.02096\ ($+0.00056$) \\
3d  & 0.70500 & 0.70526 & 0.69313 & 0.69471 & 0.01187 & 0.01056\ ($-0.00131$) \\
7d  & 0.70583 & 0.70693 & 0.69349 & 0.69660 & 0.01234 & 0.01033\ ($-0.00202$) \\
14d & 0.70667 & 0.70835 & 0.69477 & 0.69869 & 0.01190 & 0.00966\ ($-0.00223$) \\
23d & 0.70623 & 0.70862 & 0.69436 & 0.69859 & 0.01187 & 0.01003\ ($-0.00183$) \\
31d & 0.70586 & 0.70837 & 0.69440 & 0.69851 & 0.01146 & 0.00986\ ($-0.00160$) \\
\bottomrule
\end{tabular}
\end{table}

\textbf{At every $N$ from 3d through 31d, proper independent tuning
narrows, not widens, the $M_1$-vs-$M_2$ gap} (by 11--19\% of the untuned
gap), because $M_2$ gains more PR-AUC from added tree complexity than
$M_1$ does. The qualitative conclusion ($M_1 > M_2$ everywhere) is
unaffected -- the gap narrows, it does not close or reverse. At $N=2$h
and $N=1$d, tuning goes the other way (gap widens by $0.0005$--$0.0006$);
this was diagnosed as a small-sample overfitting-to-validation-split
effect specific to those two smallest budgets (a 15\% validation slice
at $N=2$h is only $\sim 1{,}080$ rows), not a competing finding, and does
not affect the $N \ge 3$d reading.

\subsection{Operational lead-time metric}
\label{sec:results-operational}

The operational lead-time metric (top 1\% alert rate, fixed budget) was
computed on the full 31-day rolling out-of-sample population
($2{,}561{,}477$ rows, prevalence $0.409426$ -- materially higher than the
fixed test week's $0.3440$, and not directly comparable to the PR-AUC/Brier
figures elsewhere in this paper). Table~\ref{tab:operational} reports the
result.

\begin{table}[htbp]
\centering
\caption{Operational lead-time metric at a fixed 1\% alert rate, 31-day
rolling OOS population (2{,}561{,}477 rows).}
\label{tab:operational}
\begin{tabular}{@{}lrrrrr@{}}
\toprule
Model & Threshold & Alert rate & Recall @ 1\% & Median lead (ms) & Mean lead (ms) \\
\midrule
$M_1$ & 0.985170 & 1.000\% & 0.024279 & 52.00 & 80.20 \\
$M_2$ & 0.965276 & 1.000\% & 0.023634 & 52.00 & 76.14 \\
$M_3$ & 0.984372 & 1.000\% & 0.024277 & 52.00 & 80.17 \\
\bottomrule
\end{tabular}
\end{table}

\textbf{$M_1$ and $M_3$ are operationally close to indistinguishable at
this budget}, on both recall ($\Delta = -0.000002$) and mean lead time
($\Delta = -0.03$\,ms), with identical median lead time (52.00\,ms) --
consistent with, not contradictory to, Section~\ref{sec:results-m3}'s
finding that $M_3$'s edge over $M_1$ is a calibration effect rather than
a top-of-ranking effect: a model that improves calibration without
improving ranking is expected to show exactly this pattern. \textbf{$M_2$
is measurably worse than $M_1$ on both dimensions here too} (recall
$-0.000645$, roughly 2.7\% fewer captured events at the same alert
budget; mean lead time $-4.06$\,ms, roughly 5\% less advance warning),
directionally consistent with every other $M_1$-vs-$M_2$ comparison in
this paper. Absolute recall is low for all three models at this budget
($\sim 2.4\%$), which is expected given the population's high ($41\%$)
target prevalence relative to a 1\%-of-samples alert budget, and should
be read as ``quality of the very top of the ranking,'' not absolute event
coverage.

\section{Limitations}
\label{sec:limitations}

\subsection{Single 38-day calendar window, single exchange}
\label{sec:limitations-window}

This is the project's own most-flagged open item and is stated
explicitly rather than left implicit. Every result in this paper
describes BTCUSDT (and, for the transfer ladder, ETH/SOL/DOGE USDT)
perpetual-futures microstructure on Binance during a single 38-day
calendar window, 2024-02-22 -- 2024-03-30. No claim in this paper should
be read as asserting that market microstructure dynamics are
time-invariant, or that the specific magnitudes reported (the
$M_1$-vs-$M_2$ PR-AUC gap, the branching-ratio values, the transfer
ladder's captured-gap percentages) generalize to other calendar periods,
volatility regimes, or exchanges without further testing. This is a
standard scope limitation for microstructure econometrics, not a defect
specific to this study, but it bears directly on external validity: the
mechanisms studied here -- order-flow clustering, cross-excitation
between quote and trade events, queue depletion -- are not expected to
have fundamentally changed in kind since early 2024, but their exact
parameterization (kernel weights, branching ratios, the $M_1$-vs-$M_2$
PR-AUC gap magnitude) should not be assumed stable across different
market regimes (e.g.\ markedly higher or lower realized volatility, a
different funding-rate environment, or a different exchange's matching
engine and tick-size conventions) without a fresh fit. The
window itself was not a free choice: it is bounded on the recent end by
Binance's permanent discontinuation of the \texttt{bookTicker} feed for
BTCUSDT futures after 2024-03-30 (Section~\ref{sec:data-feeds}), so a
more recent replication under the identical protocol is not currently
possible against this exact data source. This was re-checked directly
rather than assumed stale: as of 2026-09-18 the live archive still lists
nothing past 2024-03-30, including a monthly rollup file one step past
the daily cutoff that turned out to contain only 6h47m of 2024-04-01
data. As the nearest feasible substitute, Section~\ref{sec:results-robustness}
reruns the core learning curve against an independent, non-overlapping
test week drawn from \emph{within} the same 38-day window and finds the
same $M_1>M_2$ result -- a within-window robustness check that
corroborates the finding's stability inside this window, but does not,
and cannot, speak to external validity across a genuinely different
calendar period or market regime.

\subsection{The excluded 2024-02-04 -- 2024-02-21 disorder window}
\label{sec:limitations-disorder}

A full-coverage audit of the originally planned 60-day window found
severe raw \texttt{bookTicker} ordering disorder between 2024-02-04 and
2024-02-21 -- on some days in this range, up to $\sim$40\% of a day's
rows showed a \texttt{transaction\_time} decrease relative to the
previous row (e.g.\ 18{,}728{,}015 of 46{,}396{,}761 rows on 2024-02-12),
a scale of disorder far beyond isolated same-millisecond ties that the
project's standard defensive sort (by \texttt{transaction\_time}, then
\texttt{update\_id}) produces \emph{a} valid-looking chronological order
for, but whose recovery of the \emph{true} underlying order at this scale
of scrambling was never independently verified. This 18-day span was
excluded from the primary window entirely rather than included with an
unverified repair, and no Hawkes fit, learning-curve point, or any other
result in this paper touches it. This is a conservative, disclosed
exclusion, not a silent gap: the resulting 38-day window is the largest
\emph{contiguous, independently re-audited clean} span available between
the resolved end of this disorder period (2024-02-22, confirmed via a
zero-raw-decrease re-audit) and the permanent \texttt{bookTicker}
discontinuation (2024-03-30). Whether the excluded period's underlying
market dynamics (which, unlike the archive's row ordering, were not
inherently anomalous) would change any of this paper's findings if
recovered and independently validated is untested and remains open.

\subsection{Cross-machine Hawkes non-bit-reproducibility}
\label{sec:limitations-reproducibility}

Several results in this paper (e.g.\ Table~\ref{tab:m3-comparison}'s
$M_2$ column differing from Table~\ref{tab:learning-curve} by
$\sim 0.0001$--$0.0002$ PR-AUC at every $N$, despite an identical
specification) reflect a finding that was itself investigated and
resolved during the project: the unconstrained L-BFGS-B Hawkes fit used
throughout this paper is \emph{not} bit-reproducible across different
machines or software environments, even at an identical fixed random
seed, identical data, and identical code -- most likely due to different
BLAS/LAPACK builds shifting the optimizer's floating-point trajectory.
This was distinguished directly from true non-determinism: two
back-to-back reruns of an identical 100{,}000-event calibration budget
\emph{on the same machine} were confirmed bit-identical to full displayed
precision, while the same specification run on a different machine in an
earlier session differed by up to $0.0016$ PR-AUC -- roughly an order of
magnitude larger than the previously assumed ``run-to-run noise floor''
of $\sim0.0001$--$0.0002$, which in retrospect was itself cross-machine
noise mislabeled as generic non-determinism. LightGBM fits ($M_0$/$M_1$,
which never call the Hawkes optimizer) reproduce exactly across every
run and every machine used in this project, at every reported precision.
Every cross-model or cross-configuration \emph{comparison} reported in
this paper (e.g.\ every $\Delta$ column, every GW statistic) uses model
outputs produced within a single run on a single machine, so this
non-reproducibility does not threaten any comparison's internal validity;
it does mean that a reader attempting to reproduce an individual $M_2$ or
$M_3$ PR-AUC figure to full displayed precision on different hardware
should expect a difference on the order of $10^{-4}$--$10^{-3}$, not
bit-exact agreement, and any future rerun on a \emph{single, fixed}
machine should reproduce exactly -- a discrepancy under those specific
conditions would indicate a real bug, not expected Hawkes noise.

\subsection{Transfer-ladder simplifications, resolved}
\label{sec:limitations-transfer}

The cross-asset transfer ladder (Section~\ref{sec:results-transfer}) was
disclosed, when first built, as resting on three simplifications relative
to the specification's stated protocol: transferring the full
$10\times10\times5$ excitation tensor rather than only the collapsed
branching-ratio matrix (a strict superset of what is needed, not a
simplification against the protocol); estimating the target asset's
baseline intensity $\boldsymbol\mu$ via a closed-form empirical rate
rather than a properly profiled maximum-likelihood re-estimation under
the fixed, transferred $\boldsymbol\alpha$; and transferring BTC's
$\boldsymbol\Gamma$ directly to each of ETH, SOL, and DOGE rather than
through a literal daisy chain (each asset's own re-estimated
$\boldsymbol\Gamma$ feeding the next rung). All three were subsequently
tested directly rather than left as caveats. The literal daisy chain was
built and evaluated and produced results within $0.0008$--$0.0009$ PR-AUC
of the direct-transfer design, with no consistent direction. A proper
fixed-$\boldsymbol\alpha$, profiled-$\boldsymbol\mu$ estimator (solved by
one-dimensional bisection, including an explicit boundary-clamp case
validated against the joint MLE's own boundary solutions to within
0.08\% relative error) was implemented and used to rerun every transfer
evaluation; every resulting PR-AUC changed by at most $0.0003$, with no
consistent direction across the eight cases tested. \textbf{These
simplifications are therefore resolved, not open}: the transfer-ladder
results reported in Section~\ref{sec:results-transfer} and
Table~\ref{tab:transfer-synthesis} reflect the originally disclosed
simplified design, and this subsection records that a substantially more
faithful re-implementation of the specified protocol was subsequently
confirmed to change nothing materially. What remains genuinely untested
for the transfer ladder is orthogonal to these three points: only one
BTC source calibration window was ever used as the transfer origin, and
no formal significance test (Giacomini--White or Romano--Wolf) was run on
any transfer rung -- every transfer number in this paper is a single
train/test point estimate on its target asset, unlike the BTC-only
results elsewhere in this paper, which are corroborated by a formal
rolling test. A further robustness extension to three more assets
(BNB, XRP, ADA, Section~\ref{sec:results-transfer}) also revealed a
new, unresolved question rather than only confirming the original
finding: those three rungs cluster $1.3$ percentage points apart at
$78.97\%$--$80.27\%$ gap-capture, below the original three-asset range's
floor of $81.7\%$, and the qty-tie dual-$\boldsymbol\Gamma$ check above
was not repeated for them. Whether this reflects a structural property
of the newly added assets or is a coincidence of six single-window,
single-source point estimates is not resolved in this paper.

\subsection{Other disclosed scope decisions}
\label{sec:limitations-other}

A small number of additional scope decisions were made explicitly during
the project and are recorded here for completeness rather than treated
as open questions requiring resolution in this paper. The post-hoc
stationarity projection (Section~\ref{sec:methodology-hawkes}) was
retained rather than replaced with an in-fit constraint, on the basis of
an empirical scan across 185 fits in which the projection essentially
never bound, later found to bind in exactly one hourly window that traces
to a known, pre-existing \texttt{bookTicker} data gap
(2024-03-08) rather than to a general large-$N$ property; the affected
window contributes 0.116\% of the $N=31$d training rows and a $0.47\%$
coefficient-level rescale for that one hour. A residual $8.4\%$
tie-order sensitivity in the fitted $\boldsymbol\Gamma$ matrix, after
masking the dominant same-millisecond touch-price/trade ambiguity
(Section~\ref{sec:data-tiebreak}), was traced to a narrower, disclosed
class of same-millisecond quantity-only quote/trade ties (6.35\% of
events) that the current masking scope does not cover; this scope
decision was made explicitly and re-reviewed (and reconfirmed adequate)
before both the $M_3$ build and the transfer ladder, each of which
consumes $\boldsymbol\Gamma$ in a different way. The difference-in-differences
placebo control specified for validating $M_3$'s structural claim
(Section~\ref{sec:methodology-placebo}) was not run in this study: given
$M_3$'s null PR-AUC result on the Primary Endpoint
(Section~\ref{sec:results-m3}), it was judged to have comparatively
little to falsify at the time relative to isolating the source of $M_3$'s
subsequently discovered calibration effect via ablation, which was run
instead and is reported in Section~\ref{sec:results-m3-ablation}; running
the placebo control against the calibration effect specifically (rather
than the null PR-AUC result it was originally scoped against) remains a
natural next step. Finally, the secondary/exploratory grid's
\texttt{target\_a} implementation does not take a depletion-fraction
argument, so the nominal 18 named specifications in
Table~\ref{tab:romano-wolf} contain only 12 distinct tests; this is a
codebase quirk disclosed at the time it was discovered, not silently
corrected, and all 18 named columns are reported as specified.

\section{Conclusion}
\label{sec:conclusion}

This paper set out to test, under a falsification-first protocol, whether
a maximum-likelihood-constrained multivariate Hawkes process generalizes
better than an unconstrained gradient-boosted-tree model at forecasting
L1 order-book queue depletion, in a setting deliberately constructed so
that the parametric model's intensity is a strict linear subspace of the
tree model's own feature space -- making the comparison one of
generalization under constraint, not of which model can see more
information. Across every axis tested -- the full sample-size learning
curve from 2 hours to 31 days of training data, a formal 31-day rolling
Giacomini--White test, a calibration-budget robustness sweep up to
50$\times$ the original Hawkes calibration sample, an 18-specification
Romano--Wolf FWER-controlled grid, and independently tuned LightGBM
hyperparameters for both models -- the free tree model $M_1$ beats the
parametric Hawkes model $M_2$, by a wide and statistically decisive
margin, with no crossover sample size at which the parametric model
overtakes it. The single standing hypothesis that could have explained
this away as a methodological artifact in $M_1$'s favor -- that $M_1$'s
hyperparameters were undertuned relative to $M_2$'s -- was tested
directly and found backwards: proper tuning narrows the gap by
benefiting $M_2$ more than $M_1$, not the reverse.

At the same time, the project did not simply confirm a null and stop.
$M_3$, the one model in the ladder with access to genuinely non-linear
information ($M_1$'s raw feature bank does not already contain) -- an
hourly Hawkes-derived spectral-radius regime indicator and linearized
trade-size marks -- showed no ranking-quality (PR-AUC) advantage over
$M_1$ on the 500\,ms Primary Endpoint, exactly as the hourly regime
feature's own coarser timescale would predict. Yet a formal test on a
different, calibration-sensitive metric (log-loss, under the same
31-day rolling protocol that established the $M_1$-vs-$M_2$ result)
revealed a small but highly significant improvement, which an ablation
traced cleanly to the regime feature alone, with the trade-size marks
contributing nothing and mildly diluting the combined effect. This
result -- real on one metric, absent on another, and not large enough to
clear the pre-registered PR-AUC magnitude threshold under either reading
-- is reported as exactly that: a modest, metric-specific finding, not
rounded up into a headline reversal of the paper's primary conclusion,
consistent with the project's stated commitment to reporting what the
data says rather than what would make the best story.

The cross-asset transfer ladder adds a further, independent line of
evidence on the value of the parametric compression itself, separate
from the sample-efficiency question the learning curve addresses. A
branching-ratio matrix fit purely on BTC, transferred with only a local
baseline intensity and marks re-estimated per target, recovers 81.7\%--89.3\%
of each of three structurally different assets' (ETH, SOL, DOGE) own
free-model advantage over the $M_0$ baseline -- with no evidence of
degradation as tick size and event mix diverge further from BTC's own.
This is a genuine, if narrower, generalization advantage for the
parametric representation: not better absolute forecasting accuracy than
a locally-trained free model, but a substantial fraction of that
accuracy transported across instruments at essentially zero re-training
cost, which a black-box tree model has no comparable native mechanism to
achieve.

Taken as a whole, the evidence assembled here falsifies the hypothesis
that MLE-constrained parametric compression outperforms unconstrained
tree learning for L1 queue-depletion \emph{ranking} quality at any tested
sample size, training-data diversity, or calibration budget on this
instrument and window -- a clean, negative, and well-powered result in
its own right. It simultaneously identifies two narrower channels in
which the parametric structure retains real value: a calibration-only
predictive edge from a single Hawkes-derived regime summary, and a
substantial, non-degrading transfer of predictive structure across
instruments that a purely data-driven free model cannot replicate
without retraining. Future work most directly motivated by this paper's
own findings and limitations
(Section~\ref{sec:limitations}) includes replicating the primary
comparison over additional calendar windows and exchanges to establish
temporal and cross-venue external validity, running the specified
difference-in-differences placebo control against $M_3$'s
now-established calibration effect rather than its originally tested
null PR-AUC result, and extending the transfer ladder with a formal
significance test and additional BTC source windows now that its
disclosed simplifications have been resolved.


\bibliographystyle{plainnat}
\bibliography{../references}

\end{document}
